%%%%%%%%%%  scan idx 58  |  folio 45  %%%%%%%%%%
de dimension cohomologique finie (on a donc $\xi(X)\subset\cF_0(X)\subset\cF(X)$). Alors l'homomorphisme canonique}
\begin{equation}\tag{2.12}
 K(\xi(X))\longrightarrow K(\cF_0(X))
\end{equation}
\emph{est un isomorphisme.}

On va vérifier les conditions du théorème 2.3. Les conditions (i) et (iii) résultent de la définition de la dimension cohomologique et du fait qu'un module projectif de type fini sur l'anneau local noethérien $\cO_x$ est libre. Pour vérifier (ii), on va démontrer que $\cC=\cF(X)$ satisfait aux conditions (ii a) et (ii b) de la remarque. La deuxième est triviale~; il reste donc à prouver le résultat suivant~:

\textbf{Proposition 2.4. (Serre).} \emph{Soit $X$ un espace algébrique quasi-projectif. Alors tout faisceau algébrique cohérent sur $X$ est isomorphe au quotient d'un faisceau algébrique localement libre.}

Cette proposition est vraie si $X$ est une partie fermée de l'espace projectif en vertu de [2], puisque, pour $n$ assez grand, $F(n)$ est engendré par ses sections, ce qui signifie que $F$ est un quotient de $\cO(-n)^k$. Dans le cas général, $X$ est une partie ouverte d'une partie fermée d'un espace projectif, et il suffit d'utiliser le résultat suivant (dû à Cartier et Serre dans des cas particuliers)~:

\textbf{Proposition 2.5.} \emph{Soit $F$ un faisceau algébrique cohérent sur une partie ouverte $U$ d'un espace algébrique $X$. Alors $F$ est isomorphe à la restriction d'un faisceau algébrique cohérent sur $X$.}

En zornifiant sur les ouverts contenant $U$ sur lesquels on peut prolonger $F$, on est ramené au cas où $X$ est affine. Soit $\overline F$ le faisceau
algé%
